\documentclass[10pt,a4paper]{article}
\usepackage[a4paper,margin=14mm,top=12mm,bottom=14mm]{geometry}
\usepackage{fontspec}
\setmainfont{Liberation Sans}
\usepackage{polyglossia}
\setdefaultlanguage{french}
\usepackage{amsmath}
\usepackage{xcolor}
\usepackage{graphicx}
\usepackage[most]{tcolorbox}
\usepackage{tabularx,booktabs,array,multirow}
\usepackage{enumitem}
\usepackage{fancyhdr}
\usepackage{lastpage}
\usepackage[hidelinks]{hyperref}
\definecolor{uimm}{HTML}{005677}
\definecolor{cyan}{HTML}{13A2DD}
\definecolor{fond}{HTML}{E6F2FA}
\definecolor{ok}{HTML}{2E7D32}
\definecolor{ko}{HTML}{B3261E}
\definecolor{violet}{HTML}{7B1FA2}
\setlength{\parindent}{0pt}
\setlength{\parskip}{3pt}
\setlist{nosep}
\setlist[itemize]{leftmargin=13pt,label=\textcolor{cyan}{\textbullet}}
\newcommand{\titre}[1]{\vspace{5pt}{\color{uimm}\large\bfseries #1}\par\vspace{-3pt}{\color{cyan}\rule{\linewidth}{1.2pt}}\par\vspace{0pt}}
\tcbset{boite/.style={colback=fond,colframe=uimm,boxrule=0.6pt,arc=2pt,left=6pt,right=6pt,top=4pt,bottom=4pt},
        defi/.style={colback=violet!6,colframe=violet,boxrule=0.8pt,arc=2pt,left=6pt,right=6pt,top=4pt,bottom=4pt,
                     title={\bfseries Défi},coltitle=white,colbacktitle=violet,fonttitle=\small}}
\newcommand{\bandeau}[3]{%
\begin{tcolorbox}[colback=uimm,colframe=uimm,arc=2pt,left=4pt,right=8pt,top=4pt,bottom=4pt]
\begin{minipage}[c]{18mm}\setlength{\fboxsep}{1.2pt}\fcolorbox{white}{white}{\includegraphics[height=16.5mm]{logo_uimm.jpg}}\end{minipage}\hspace{4mm}%
\begin{minipage}[c]{\dimexpr\linewidth-22mm\relax}\color{white}
{\small #1}\\[1pt]{\LARGE\bfseries #2}\\[2pt]{\large #3}\hfill{\small CQPM Technicien qualité}
\end{minipage}
\end{tcolorbox}}
\newcommand{\reponse}[1][1]{\par\vspace{2pt}{\color{gray!60}\foreach \i in {1,...,#1}{\rule{\linewidth}{0.3pt}\par\vspace{9pt}}}}
\usepackage{pgffor}
\pagestyle{fancy}\fancyhf{}
\renewcommand{\headrulewidth}{0pt}
\newcommand{\pied}[1]{\fancyfoot[L]{\footnotesize\color{uimm}#1}\fancyfoot[R]{\footnotesize\color{uimm}Page \thepage{} sur \pageref{LastPage}}}

\usepackage{tikz}
\usepackage{longtable}
\renewcommand{\bandeau}[3]{%
\begin{tcolorbox}[colback=uimm,colframe=uimm,arc=2pt,left=4pt,right=8pt,top=4pt,bottom=4pt]
\begin{minipage}[c]{18mm}\setlength{\fboxsep}{1.2pt}\fcolorbox{white}{white}{\includegraphics[height=16.5mm]{logo_uimm.jpg}}\end{minipage}\hspace{4mm}%
\begin{minipage}[c]{\dimexpr\linewidth-22mm\relax}\color{white}
{\small #1}\\[1pt]{\LARGE\bfseries #2}\\[2pt]{\large #3}
\end{minipage}
\end{tcolorbox}}
% en-tête d'une fiche : code, titre, question, bloc, durée, organisation, niveau
\newcommand{\entete}[7]{%
\begin{tcolorbox}[colback=uimm,colframe=uimm,arc=2pt,left=4pt,right=6pt,top=4pt,bottom=4pt]
\begin{minipage}[c]{18mm}\setlength{\fboxsep}{1.2pt}\fcolorbox{white}{white}{\includegraphics[height=16.5mm]{logo_uimm.jpg}}\end{minipage}\hspace{3mm}%
\begin{minipage}[c]{\dimexpr\linewidth-58mm\relax}\color{white}\raggedright
{\small Base SPC · #4}\\[1pt]{\Large\bfseries #1 · #2}\\[2pt]{\itshape #3}
\end{minipage}\hfill
\begin{minipage}[c]{33mm}\color{white}\small\raggedleft
Durée : \textbf{#5}\\ \textbf{#6}\\ Niveau #7 sur 3
\end{minipage}
\end{tcolorbox}}
\newcommand{\tagE}{{\setlength{\fboxsep}{1.5pt}\colorbox{cyan}{\color{white}\scriptsize\bfseries Ellistat}}}
\newcommand{\tagW}{{\setlength{\fboxsep}{1.5pt}\colorbox{uimm}{\color{white}\scriptsize\bfseries Page web}}}
\tcbset{pred/.style={colback=orange!6,colframe=orange!80!black,boxrule=0.6pt,arc=2pt,left=6pt,right=6pt,top=3pt,bottom=1pt,
                    title={\bfseries Je prédis, avant de commencer},coltitle=white,colbacktitle=orange!80!black,fonttitle=\small},
        retenir/.style={colback=fond,colframe=uimm,boxrule=0.6pt,arc=2pt,left=6pt,right=6pt,top=3pt,bottom=3pt,
                    title={\bfseries À retenir},coltitle=white,colbacktitle=uimm,fonttitle=\small},
        plus/.style={colback=violet!6,colframe=violet,boxrule=0.6pt,arc=2pt,left=6pt,right=6pt,top=3pt,bottom=3pt,
                    title={\bfseries Pour aller plus loin},coltitle=white,colbacktitle=violet,fonttitle=\small},
        act/.style={colback=uimm!10,colframe=uimm,boxrule=0.5pt,arc=2pt,left=5pt,right=5pt,top=2pt,bottom=2pt,before skip=7pt,after skip=3pt},
        verif/.style={colback=orange!6,colframe=orange!80!black,boxrule=0.6pt,arc=2pt,left=6pt,right=6pt,top=3pt,bottom=3pt,
                    title={\bfseries À vérifier sur poste avant la première séance},coltitle=white,colbacktitle=orange!80!black,fonttitle=\small}}
\newcommand{\fichier}[1]{\texttt{\small #1}}
\newcommand{\menu}[1]{\textsf{\textbf{#1}}}
\newcommand{\champ}[1]{\rule{#1}{0.3pt}}

\pied{Base SPC · fiche C4 · Ma carte, à la main puis dans Ellistat}
\begin{document}
\thispagestyle{fancy}
\entete{C4}{Ma carte, à la main puis dans Ellistat}{La carte d'Ellistat dit-elle la même chose que la mienne ?}{Cartes aux mesures et capabilité}{60 min}{en binôme}{2}
\begin{tcolorbox}[boite]\textbf{Le but.} Calculer les limites d'une carte des moyennes et des étendues, puis les comparer à celles d'Ellistat.\par\textbf{Il vous faut.} La grille des 30 butées en 6 sous-groupes (activité E5) · Une calculatrice\end{tcolorbox}
\begin{tcolorbox}[pred]Un logiciel et un calcul à la main peuvent-ils donner des limites légèrement différentes et avoir raison tous les deux ?\reponse\end{tcolorbox}
\titre{À faire}
\begin{enumerate}[label=\textcolor{cyan}{\textbf{\arabic*.}},leftmargin=15pt,itemsep=1.5pt]
\item Calculer la moyenne et l'étendue de chacun des 6 sous-groupes, puis $\bar{\bar X}$ et $\bar R$.
\item Limites : $LCS = \bar{\bar X} + 0{,}577\,\bar R$ ; $LCI = \bar{\bar X} - 0{,}577\,\bar R$ ; pour les étendues $LCS_R = 2{,}114\,\bar R$.
\item Écart-type court terme : $\bar R / 2{,}326$. Puis $Cp = 0{,}15 / (6 \times \text{écart-type})$.
\item \tagE\ \menu{Stats Desc.} > \menu{Détail}, référentiel \menu{Centré Moyenne}. Relever les limites, l'écart-type court terme, Cp, Cpk, Pp, Ppk.
\end{enumerate}
\vspace{3pt}{\renewcommand{\arraystretch}{1.45}\begin{tabularx}{\linewidth}{|l|*{6}{>{\centering\arraybackslash}X|}}\hline
\textbf{Sous-groupe} & 1 & 2 & 3 & 4 & 5 & 6\\\hline Moyenne & & & & & & \\\hline Étendue $R$ & & & & & & \\\hline\end{tabularx}}\par\vspace{3pt}
{\renewcommand{\arraystretch}{1.4}\begin{tabularx}{\linewidth}{|X|c|c|}\hline & \textbf{À la main} & \textbf{Ellistat}\\\hline
Moyenne générale $\bar{\bar X}$ \quad / \quad étendue moyenne $\bar R$ & \hspace{28mm} & \hspace{28mm}\\\hline
Limite supérieure des moyennes & \hspace{28mm} & \hspace{28mm}\\\hline
Limite inférieure des moyennes & \hspace{28mm} & \hspace{28mm}\\\hline
Limite supérieure des étendues & \hspace{28mm} & \hspace{28mm}\\\hline
Écart-type court terme & \hspace{28mm} & \hspace{28mm}\\\hline
Cp \quad / \quad Cpk & \hspace{28mm} & \hspace{28mm}\\\hline
Pp \quad / \quad Ppk & \hspace{28mm} & \hspace{28mm}\\\hline
\end{tabularx}}\par
\titre{J'observe, je décide}
\begin{enumerate}[label=\textcolor{cyan}{\textbf{Q\arabic*.}},leftmargin=22pt]
\item Les mêmes points sortent-ils des limites dans les deux calculs ?\reponse[1]
\item Ellistat utilise toutes les mesures de chaque sous-groupe, l'étendue seulement deux. Lequel est le plus précis ?\reponse[1]
\item Cp est correct mais Cpk est faible. Que montre l'histogramme ? Passez en \menu{Centré Cible}.\reponse[2]
\item Pp est plus petit que Cp. Que s'est-il passé entre les sous-groupes ?\reponse[1]
\end{enumerate}
\begin{tcolorbox}[retenir]Deux méthodes, deux chiffres voisins, une même décision. Cp regarde l'intérieur des sous-groupes, Pp toute la période.\end{tcolorbox}
\vfill{\footnotesize\color{uimm}\textbf{Ressources de la base :} Données à copier.}
\end{document}